\begin{trapbox}
	\begin{description}[style=nextline, leftmargin=2.8em, labelsep=0.5em, itemsep=0.35em]
		\item[\textbf{\textcolor{CTrap}{易错点一（遗漏分母条件）}}]
		      当 $\lambda=1$ 时，若仍直接使用
		      \[
			      C\left(\frac{\lambda^2+1}{\lambda^2-1}a,\;0\right),
			      \qquad
			      R=\left|\frac{2a\lambda}{\lambda^2-1}\right|,
		      \]
		      就会出现分母 $\lambda^2-1=0$ 的问题.

		\item[\textbf{\textcolor{CTrap}{正确理解（分类讨论）：}}]
		      当 $\lambda=1$ 时，条件
		      \[
			      \frac{PA}{PB}=1
		      \]
		      表示 $P$ 到 $A,B$ 的距离相等，轨迹是线段 $AB$ 的中垂线；只有当 $\lambda>0$ 且 $\lambda\neq1$ 时，轨迹才是阿波罗尼斯圆.

		\item[\textbf{\textcolor{CTrap}{错因分析（公式适用范围）：}}]
		      学生直接套公式，忽略分母不能为零，也忽略了 $\lambda=1$ 时轨迹类型已经改变.
	\end{description}

	\medskip

	\begin{description}[style=nextline, leftmargin=2.8em, labelsep=0.5em, itemsep=0.35em]
		\item[\textbf{\textcolor{CTrap}{易错点二（半径符号处理）}}]
		      直接写
		      \[
			      R=\frac{2a\lambda}{\lambda^2-1}
		      \]
		      而不加绝对值. 当 $0<\lambda<1$ 时，分母 $\lambda^2-1<0$，这样会得到负数，不能作为半径.

		\item[\textbf{\textcolor{CTrap}{正确理解（绝对值取正）：}}]
		      由于 $a>0,\lambda>0,\lambda\neq1$，所以
		      \[
			      R=\left|\frac{2a\lambda}{\lambda^2-1}\right|>0.
		      \]
		      半径必须是严格正值，因此表达式外必须加绝对值.

		\item[\textbf{\textcolor{CTrap}{错因分析（开平方处理）：}}]
		      从圆方程右边开方时，半径取的是算术平方根，不能保留可能为负的代数式.
	\end{description}

	\medskip

	\begin{description}[style=nextline, leftmargin=2.8em, labelsep=0.5em, itemsep=0.35em]
		\item[\textbf{\textcolor{CTrap}{易错点三（坐标不对称）}}]
		      当题目给出的两定点不是 $(-a,0)$ 和 $(a,0)$ 的形式时，学生机械代入标准公式，容易把圆心和半径算错.

		\item[\textbf{\textcolor{CTrap}{正确理解（一般坐标版）：}}]
		      若 $A(x_1,y_1)$，$B(x_2,y_2)$，且
		      \[
			      \frac{PA}{PB}=\lambda,\qquad \lambda>0,\quad \lambda\neq1,
		      \]
		      应先写
		      \[
			      (x-x_1)^2+(y-y_1)^2
			      =
			      \lambda^2\bigl[(x-x_2)^2+(y-y_2)^2\bigr].
		      \]
		      整理后圆心和半径为
		      \[
			      C\left(
			      \frac{x_1-\lambda^2x_2}{1-\lambda^2},
			      \frac{y_1-\lambda^2y_2}{1-\lambda^2}
			      \right),
			      \qquad
			      R=
			      \frac{\lambda\sqrt{(x_1-x_2)^2+(y_1-y_2)^2}}{|1-\lambda^2|}.
		      \]

		\item[\textbf{\textcolor{CTrap}{错因分析（死记硬背）：}}]
		      学生只记忆标准坐标公式，不理解公式来自距离方程的平方和配方，因此一遇到非对称坐标就容易误套.
	\end{description}

	\medskip

	\begin{description}[style=nextline, leftmargin=2.8em, labelsep=0.5em, itemsep=0.35em]
		\item[\textbf{\textcolor{CTrap}{易错点四（空间问题误用平面结论）}}]
		      在立体几何中，若 $P$ 是空间中的任意动点，满足
		      \[
			      \frac{PA}{PB}=\lambda,\qquad \lambda>0,\quad \lambda\neq1,
		      \]
		      其轨迹一般不是圆，而是阿波罗尼斯球面.

		\item[\textbf{\textcolor{CTrap}{正确理解（平面与空间要分清）：}}]
		      若点 $P$ 被限制在包含 $A,B$ 的某一平面内运动，轨迹是该平面内的阿波罗尼斯圆；若 $P$ 在空间中自由运动，轨迹是阿波罗尼斯球面. 若再与某个平面相交，交线才可能是圆、一个点或空集.

		\item[\textbf{\textcolor{CTrap}{错因分析（维度混淆）：}}]
		      平面内的距离比轨迹是圆，空间内的距离比轨迹是球面；本质相同，但维度不同，不能直接把空间轨迹说成圆.
	\end{description}
\end{trapbox}